% ECONOMICS_PROBLEM: {"id": "D63-84", "title": "Linear total subsidies for general monotone goods valuations", "jel_code": "D63", "source_id": "OP-0591"}
% !TeX program = xelatex
\documentclass[11pt,a4paper]{article}
\usepackage[margin=23mm]{geometry}
\usepackage{fontspec}
\setmainfont{Latin Modern Roman}
\usepackage{amsmath,amssymb,mathtools}
\usepackage{xcolor,enumitem,booktabs,longtable,array,xurl,hyperref}
\definecolor{muted}{HTML}{53636C}
\hypersetup{colorlinks=true,urlcolor=blue,linkcolor=blue}
\setlength{\parindent}{0pt}
\setlength{\parskip}{5pt}
\newcommand{\R}{\mathbb R}
\newcommand{\E}{\mathbb E}
\newcommand{\N}{\mathbb N}
\newcommand{\ind}{\mathbf 1}
\newcommand{\Var}{\operatorname{Var}}
\newcommand{\Cov}{\operatorname{Cov}}
\newcommand{\supp}{\operatorname{supp}}
\DeclareMathOperator*{\argmin}{arg\,min}
\DeclareMathOperator*{\argmax}{arg\,max}
\newcommand{\entrymeta}[3]{{\sffamily\small\color{muted}\textbf{\texttt{#1}}\quad|\quad #2\par #3\par}}
\newcommand{\Problemref}[1]{\texttt{#1}}
\newcommand{\JELref}[2]{\texttt{#2} (JEL \texttt{#1})}
\begin{document}
\section*{Linear total subsidies for general monotone goods valuations}
\textbf{Problem ID:} \texttt{D63-84}\quad \textbf{JEL:} \texttt{D63}\par
% BEGIN ECONOMICS STATEMENT
\label{problem:OP-0591}
% ECONOMICS_PROBLEM: {"local_id": "ECON-004", "title": "Linear total subsidies for general monotone goods valuations", "kind": "U", "audit_date": "2026-10-08", "source_identity_checked": true, "agenda_or_question_support_checked": false, "source_role": "see per-entry audit and locators", "current_open_certified": false}


\label{audited:ECON-004}

{\sffamily\small\color{muted}\textbf{ECON-004} | U | Source passage unverified; not certified as a published open question\par}

\subsection*{Definitions and formal target}

Let $v_i:2^M\to\mathbb R_{\geq0}$ be normalized and monotone. Assume $0\leq v_i(S\cup\{g\})-v_i(S)\leq1$ for $g\notin S$. Subsidies $p_i\geq0$ yield utility $v_i(A_j)+p_j$ to agent $i$ from receiving bundle-payment pair $j$. Define

\[
 s(v)=\inf_{A,p\geq0}\left\{\sum_i p_i:
 v_i(A_i)+p_i\geq v_i(A_j)+p_j\quad\forall i,j\right\},
\]
where $A$ ranges over complete partitions. Is there a universal constant $C<\infty$ such that $s(v)\leq Cn$ for every $n,M,v$ in this class?

\subsection*{Clarifications and required assumptions}

Normalized means $v_i(\varnothing)=0$; monotone means $v_i(S)\leq v_i(T)$ whenever $S\subseteq T$. The marginal bound holds for every $S\subseteq M$ and $g\notin S$, not merely for singleton bundles. Use an infimum in the definition of $s(v)$, with $\inf\varnothing=+\infty$, if attainment or feasibility has not separately been established. The constant $C$ must be independent of $n,m$ and the valuation profile. The subsidy travels with the compared bundle: the right side contains $p_j$, not $p_i$.

\subsection*{Variants and changes of scope}

\textbf{Scaling.} A marginal bound $B>0$ changes a linear target to $CBn$. \textbf{Broader sign domain.} Nonnegative bundle values without monotonicity are a larger class than monotone goods. \textbf{Transfers.} Allowing negative payments or requiring $\sum_i p_i=0$ creates a budget-balanced transfer problem, rather than a nonnegative subsidy problem. \textbf{Algorithmic version.} Bound oracle queries and bit operations explicitly; the existence quantifier alone contains no algorithm. None of these variants is attributed as explicitly open to the unavailable full text.

\subsection*{Audit conclusion and source relationship}

\textbf{Verification incomplete.} The paper's title, authors, arXiv identifier, and broad subsidy topic were corroborated, but its full text was unavailable in this audit. The earlier precise Section 6 attribution has been removed. This entry is a rigorously specified candidate target; its exact attribution and current open status are not certified. It must not be counted as a verified published open problem.

\subsection*{References and exact source roles}

{\footnotesize\raggedright
\par\noindent\textbf{[S1]} Max Dupr\'e la Tour and Mashbat Suzuki (2026). \emph{Subquadratic Subsidies for Nonnegative or Nonpositive Valuations}. arXiv:2609.08272v1, 8 September 2026.\par \textit{Verified locator / source role:} Bibliographic identity and broad abstract topic corroborated by the indexed arXiv record. Full text could not be retrieved. The draft\textquotesingle s claimed Section 6 locators are NOT verified and have been withdrawn.\par \url{https://arxiv.org/abs/2609.08272}\par\smallskip
}

\subsection*{Common conventions: Additional-source status and mathematical conventions}

\label{audited:conventions}

Of the 100 supplied ECON entries, 65 distinct problems appear here; 32 close
matches refine earlier entries, and three duplicate questions map to existing
entries. Appendix~\ref{app:audited-crosswalk} lists every destination.
The attachment's P/R/S/U distinctions and source audits are retained. Its
precise mathematical formalizations are editorial unless the individual audit
says otherwise. Candidate subsidy targets ECON-004--006 retain their U flags;
the resolved approval-core question ECON-007 updates OP-0202 above.
The following supplied conventions govern both this part and the attributed
ECON refinements elsewhere in the catalogue, subject to local restrictions.

\textbf{Parameterized questions.} A problem family lists inputs that must be fixed before an instance is evaluated: population, horizon, policies, information, data law, model class and welfare weights. These are required choices, not quantities the citations are assumed to specify. Undefined applications should not be treated as identified estimands.

\textbf{Indices and integrability.} Sets of agents, firms and regions are finite unless a population measure is explicitly used. \(i,j\) index units, \(r\) regions, \(t\) dates and \(h\) horizons. \(\mathbb E\) and \(\Pr\) refer to the declared population and law; \(\mathbf1\{A\}\) is an indicator. Every displayed expectation and discounted sum needs existence in the stated domain. Infinite horizons use a discount in \((0,1)\) and a convergence condition; finite horizons may use discount one.

\textbf{Optimization and existence.} A displayed \(\max\), \(\min\), or \(\arg\max\) is conditional on attainment. For finite feasible menus this is immediate; general spaces need continuity and compactness/coercivity or another existence argument. Without attainment, the target is the supremum/infimum and arbitrarily near-optimal feasible choices. \(\inf\varnothing=+\infty\). Empty feasibility and an unidentified target are different issues.

\textbf{Potential outcomes.} \(Y_i(z)\) is the outcome under a specified intervention, not the observed conditional mean among people choosing \(z\). In binary comparisons \(\Delta Y_i=Y_i(1)-Y_i(0)\). Specify assignment, treatment versions, compliance, information and observation horizon. Interference is allowed under a full policy vector; compressed exposure notation needs a justified mapping. No interference, consistency and overlap are substantive assumptions, not automatic consequences of notation.

\textbf{Populations and observation.} Fix baseline eligibility and population weights. Conditioning on post-policy employment, survival, relocation or program uptake can change the population being compared. Death is not ordinary loss to follow-up; outcomes truncated by death need a composite convention, joint survival target, or explicitly defined survivor stratum. Fertility-changing interventions need a population functional rather than an unexplained fixed descendant cohort. Measurement and missingness models must be supplied.

\textbf{Identification.} A structural model \(m\in\mathcal M\) implies observable law \(P_m\) and target \(\theta(m)\). Identification means
\[
 P_{m_1}=P_{m_2}\ \Longrightarrow\ \theta(m_1)=\theta(m_2)
 \quad\text{for all }m_1,m_2\in\mathcal M .
\]
Without identification, the target is
\[
 \Theta(P;\mathcal M)=\{\theta(m):m\in\mathcal M,\ P_m=P\}.
\]
Writing \(\theta(P)\) before establishing identification can hide incompatible latent models. Confidence sets additionally need a sampling law, confidence level, asymptotic sequence and uniformity class. Point identification, coverage and informative precision are separate properties.

\textbf{Mediation.} Total effects of randomized assignment are distinct from cross-world natural indirect effects. Belief/effort-channel effects require a meaningful mediator intervention and additional measurement, exclusion or mediation assumptions. Randomization of the initial policy alone does not supply those assumptions. Report bounds or interventional alternatives when the stronger target is unsupported.

\textbf{Equilibrium.} A counterfactual \(W(z)\), \(p(z)\), or \(\mu_t^z\) requires individual optimization, feasibility, government/resource budgets, market clearing and state-transition laws. State equilibrium existence and selection, or report ranges over compatible equilibria. Initial-state integration includes subsequent endogenous transitions; current-state integration must use the policy-dependent distribution. Reduced-form invariance after policy is not assumed.

\textbf{Welfare accounts.} Scores, utility, health, dollars and ecological quantities require explicit conversion before addition. Fees, taxes, subsidies, wages and extortion payments are transfers between parties; distributional weights can make them welfare relevant, but their face value is not automatically a resource loss. State ownership, geographic boundaries, foreign-income weights and fiscal finance. Do not add profits to household welfare already including dividends, or count land capitalization and the underlying benefit twice. A benefit-cost expression must distinguish gross benefits from net welfare and subtract every real cost exactly once.

\textbf{Derivatives and risk.} Log derivatives require positive arguments; local responses require admissible perturbations and specified equilibrium branches. Differentiation under expectations needs regularity/domination, and boundaries may allow only one-sided derivatives. Risk uses a specified law; ambiguity uses a declared nonempty law set on a common history space. Adaptive policies must be nonanticipative. A robust criterion is an editorial choice unless attributed explicitly.

\textbf{Scope overlaps.} Allocation, collective choice, contracts, household bargaining and coordinated policy overlap economics with optimization and game theory. They are retained because they appeared in the original catalogue; no claim of a strictly game-theory-free selection is made.
% END ECONOMICS STATEMENT
\end{document}
